\chapter{When a statistical difference changes a banking decision}
\label{scisyn:evidence}

The case for a shared analytical service is strongest when it prevents a
material mistake, exposes a valuable uncertainty or makes a defensible study
less costly. The following examples examine those propositions directly.
They are project-generated calculations, not measurements of bank returns or
claims of superiority over an installed commercial system. A fixed design
specifies the scenarios and comparisons before the simulation is run.

\section{A precise answer to the wrong offer question}
Consider an illustrative portfolio with equal proportions of two observed
customer strata, $S=0$ and $S=1$. An offer has no effect in the first stratum
and increases the outcome by 100 USD in the second. Potential outcomes are
\[
 Y(a)=50S+100Sa+\varepsilon,\qquad \varepsilon\sim N(0,50^2).
\]
Customers are independent. The offer is assigned using only $S$, with probability
$e(0)=1/2$ and $e(1)\in\{1/2,1/10,1/50\}$. The assignment is independent of
the potential-outcome noise conditional on $S$. The population average benefit
is 50 USD. The stipulated cost of a universal offer is 40 USD, leaving 10 USD
per eligible customer. These monetary values are illustrative assumptions.

The first estimator regresses residual outcomes on residual treatment using
the exact conditional means. This favourable assumption removes nuisance-model
error so that the comparison isolates the final constant-effect restriction.
The second estimates a treatment contrast separately within each stratum and
averages the two contrasts with the known population shares. Both use the same
observations. The residual estimate uses a heteroskedastic sandwich variance;
the stratified estimate uses independent-cell sample variances. Each interval
uses a normal critical value. A stratified estimate is unavailable if any of
its four cells has fewer than two observations.

The residual projection targets
\[
 \theta_*(e)=\frac{100e(1-e)}{1/4+e(1-e)}.
\]
It equals 50 under equal assignment probabilities, $450/17\simeq26.47$ when
$e=1/10$, and about 7.27 when $e=1/50$. The second estimator continues to target
50 whenever its contrasts are estimable. The projection formula follows from
the residual-regression objective inspected in the EconML comparison; the
experiment here is an independently written reference calculation, not an
EconML execution or a claim about every estimator in that package.
\citep{exec19_econml_source}

\begin{figure}[htbp]
\centering
\begin{tikzpicture}[x=17cm,y=.065cm,>=Stealth,font=\small]
\draw[->] (0,0)--(.54,0) node[right]{$e(1)$};
\draw[->] (0,0)--(0,62) node[above]{USD};
\draw[dashed,ink] (0,50)--(.5,50);
\draw[dotted,gray] (0,40)--(.5,40);
\draw[thick,teal,domain=.001:.5,samples=70] plot (\x,{100*\x*(1-\x)/(.25+\x*(1-\x))});
\foreach \x in {.1,.2,.3,.4,.5}{\draw (\x,0)--(\x,-1);\node[below] at (\x,-1){\x};}
\node[left] at (0,50){50};\node[left] at (0,40){40};
\node[above] at (.34,51){Population benefit};
\node[below] at (.36,39){Stipulated offer cost};
\node[above,align=center] at (.15,16){Constant residual\\projection};
\end{tikzpicture}
\caption{A restricted projection changes with the assignment probabilities
even though the population's causal benefit remains 50. Exact constructed
example with equally likely strata and effects 0 and 100.}
\label{scisyn:projection-figure}
\end{figure}

There are 500 replications for each of six scenarios, crossing sample sizes
500 and 2,000 with the three assignment probabilities. The fixed seed is
20261029. Table~\ref{scisyn:offer-results} reports the observed results.
The maximum Monte Carlo standard error of any displayed conditional proportion
is about 0.023; a proportion near 0.95 has a standard error of about 0.010 when
all 500 replications are available. A zero observed failure frequency does not
prove that the true probability is zero.

\begin{table}[htbp]
\centering\small
\begin{tabular}{rrlrrrr}
\toprule
$n$ & $e(1)$ & Method & Mean & ATE cover & Reject & Fail \\
\midrule
500 & 0.50 & Constant projection & 50.3 & 0.944 & 0.030 & 0 \\
500 & 0.50 & Stratum standardisation & 50.3 & 0.942 & 0.018 & 0 \\
500 & 0.10 & Constant projection & 26.3 & 0.044 & 0.988 & 0 \\
500 & 0.10 & Stratum standardisation & 49.8 & 0.940 & 0.070 & 0 \\
500 & 0.02 & Constant projection & 7.5 & 0.000 & 1.000 & 0 \\
500 & 0.02 & Stratum standardisation & 49.8 & 0.892 & 0.223 & 11 \\
2000 & 0.50 & Constant projection & 50.1 & 0.932 & 0.000 & 0 \\
2000 & 0.50 & Stratum standardisation & 50.0 & 0.954 & 0.000 & 0 \\
2000 & 0.10 & Constant projection & 26.4 & 0.000 & 1.000 & 0 \\
2000 & 0.10 & Stratum standardisation & 50.0 & 0.962 & 0.000 & 0 \\
2000 & 0.02 & Constant projection & 7.3 & 0.000 & 1.000 & 0 \\
2000 & 0.02 & Stratum standardisation & 50.0 & 0.948 & 0.056 & 0 \\
\bottomrule
\end{tabular}
\caption{Project-generated offer study. True population benefit is 50 and
stipulated cost is 40 illustrative USD per customer. Coverage and rejection
frequencies condition on an available estimate; Fail counts unavailable
replications out of 500. These are independent reference calculations,
not outputs from an external causal package.}
\label{scisyn:offer-results}
\end{table}


At $n=2{,}000$ and $e(1)=0.1$, the mean projection is 26.447. Its interval covers
its own target in 97.2 percent of replications, but covers the population ATE in
none. A decision rule that wrongly treats it as the population benefit rejects
the profitable universal offer in all 500 replications. Stratum standardisation
averages 50.011, covers the ATE in 96.2 percent and makes no such point-rule
rejections in that scenario. Increasing precision around a different estimand
does not repair the economic interpretation.

The second estimator is not an unconditional winner. At $n=500$ and $e(1)=0.02$,
it is unavailable in 11 of 500 replications. Among the remaining 489, its normal
interval covers the ATE only 89.2 percent of the time, and its point estimate
falls below cost in 22.3 percent. Counting unavailable intervals as noncoverage
gives 87.2 percent. This is a useful adverse result: specifying the right target
does not guarantee accurate small-sample inference when treatment support is
weak. The proposal should report both the target and the information available
for estimating it, including failures.

The practical recommendation is narrow. Before comparing packages, require
the analysis to state whether it reports a population average, a conditional
effect or a projection. Compare the simplest estimator answering that question.
Where a subgroup is rarely treated, consider a narrower eligible population,
a better-designed pilot or an explicitly more conservative uncertainty method.
Those changes alter the proposed study and should be explained to its sponsor.

\section{A remedy can exchange precision for assumptions}
The adverse small-cell finding motivates a second, separately specified
experiment. It uses new random draws with $n=500$ and $e(1)=0.02$, 3,000
replications per scenario and seed 2026102902. In the first scenario all four
treatment-by-stratum cells have outcome standard deviation 50. In the second,
the rarely treated $S=1,A=1$ cell has standard deviation 150 while the other
three remain at 50. Assignment is still independent of potential-outcome noise
within stratum, and the ATE remains 50. The second scenario tests an assumption
used by one proposed remedy; it does not change the target.

Three intervals use the same stratum-standardised point estimate. The first is
the original normal Wald interval. The second pools the four within-cell
variances and uses a Student critical value with $n-4$ degrees of freedom.
Conditional on the counts, independent normal observations with a common
variance justify that pooled interval. The third constructs a two-sided Student
interval with error probability $0.05/4$ for each cell mean and combines the
four intervals with the contrast weights $(-1/2,1/2,-1/2,1/2)$. If $r_j$ is
the half-width for cell $j$, the half-width for the contrast is
$\sum_{j=1}^4 |w_j|r_j$. A union bound gives coverage of at least 95 percent
conditional on the counts under independent normal observations, allowing
different variances across cells. It is deliberately conservative.

\begin{table}[htbp]\centering\small
\begin{tabular}{llrrr}
\toprule
Noise scenario & Interval & Coverage & Mean width & Recommend offer \\
\midrule
Equal variance & Normal Wald & 0.876 & 44.7 & 0.195 \\
Equal variance & Pooled Student & 0.947 & 48.1 & 0.140 \\
Equal variance & Cellwise Student & 0.999 & 250.3 & 0.003 \\
Unequal variance & Normal Wald & 0.868 & 126.3 & 0.105 \\
Unequal variance & Pooled Student & 0.522 & 49.3 & 0.345 \\
Unequal variance & Cellwise Student & 0.995 & 629.6 & 0.005 \\
\bottomrule
\end{tabular}
\caption{New generated study of sparse-support remedies. All proportions
condition on four cells with at least two observations. There are 124
unavailable cases under equal variance and 116 under unequal variance, each
out of 3,000. Width is in illustrative USD. Recommend offer means that the
interval's lower endpoint exceeds the cost of 40.}
\label{scisyn:interval-remedies}
\end{table}

Pooling repairs much of the first scenario's undercoverage, but fails badly
when the rarely treated cell has greater variability: conditional coverage
falls to 52.2 percent. The cellwise construction retains coverage in both
scenarios at the expense of very wide intervals. Its median widths are 117.7
and 288.1 USD; a few extremely small cells drive still larger mean widths.
It almost never supports the offer through a lower bound above cost. That
hesitation is informative: the data contain little precision about the rare
treated subgroup. Unavailable replications remain unavailable under all three
methods; a conditional coverage guarantee cannot manufacture missing cells.

The comparison rules out a simple promise that an alternative critical value
will solve weak support. One remedy needs a substantively defensible common
variance; another accepts less decision-making precision under a broader
normal model. Non-normal monetary outcomes, dependent customers or estimated
population shares need additional analysis. A better-balanced experiment or a
larger sample may be the valuable next expenditure. The maximum Monte Carlo
standard error for the conditional proportions is approximately 0.0094, which
does not explain the large coverage losses in the adverse scenario.

\section{Graph ambiguity has an economic price}
The inherited pcalg example specifies a Gaussian covariance for a rate change
$X$, an intermediate response $M$ and a monetary outcome $Y$:
\[
\Sigma=\begin{pmatrix}
1&0.6&0.48\\0.6&1.36&1.088\\0.48&1.088&1.8704
\end{pmatrix}.
\]
Three noncollider chains reproduce this covariance with positive independent
error variances. Their effects of $X$ on $Y$ are $0.48,0,0$. The new calculation
reconstructs all three structural systems and checks the equality of their
observational covariances. A program can therefore reproduce the observations
while leaving a consequential intervention ambiguity. The IDA source supplies
the effect-set interpretation; the displayed system is a project construction.
\citep{exec13_pcalg_reuse_primary}

If $X$ is in percentage points and $Y$ in thousands of currency units, a
25-basis-point action has gross benefit either 0 or 120 per modelled unit.
At an illustrative implementation cost of 60, the net values are $-60$ and
$60$. There is no sign-robust recommendation across those models. A diagram
assistant that selects the forward chain without evidence would make the
decision look settled while adding no identifying information. A justified
temporal restriction or experiment could resolve the choice. No expected value
of that information is calculated without probabilities over the competing
models and an account of the experiment's precision and cost.

\section{Panel weights express different investment questions}
In the inherited two-cohort panel, early adopters have effects 2 and 4, late
adopters have an observed effect of 10, and their population shares are $3/4$
and $1/4$. Recalculation confirms averages 4, $19/4$, 4 and $15/4$ under four
aggregation rules. They are not four estimates of an identical target. The
did source's distinct aggregation branches make this an implementation-relevant
comparison.
\citep{exec24_panelr_did}

A sponsor interested in benefits per exposed branch-period needs different
weights from one comparing the typical adoption cohort or the typical calendar
period. Costs must use the same units. A table of benefits in dollars does not
make mismatched denominators comparable. The next panel experiment should
freeze the target weights before it compares estimators, retain missing cells
and test sensitivity to deviations from the maintained trend restrictions.

\section{What better text measurement would have to accomplish}
Suppose an assistant extracts distress from customer notes. Let true distress
be $U$, measured distress $W=U+\varepsilon_W$, action $A=U+\varepsilon_A$ and
outcome $Y=\tau A+U+\varepsilon_Y$. All disturbances are mutually independent,
centred Gaussian variables, with unit variance except
$\operatorname{Var}(\varepsilon_W)=s$. The coefficient on $A$ in a population
regression adjusting for $W$ is
\[
 \beta_{A\mid W}=\tau+\frac{s}{1+2s}.
\]
The calculation was independently checked by solving the covariance normal
equations. For $\tau=-1/5$, the adjusted coefficient is zero at $s=1/3$ and
positive above that threshold. This is residual confounding from imperfect
measurement, even with an arbitrarily large sample. The retained review of
text-based adjustment motivates the measurement question; the equation is
a project derivation.
\citep{exec13_census_text_review_arxiv}

The useful contribution of a language model would be to improve measurement
enough to alter the effect or decision relative to a credible alternative.
It is insufficient to show that its summary is fluent. Compare human coding,
a dictionary, conventional supervised extraction and the language model on
the same permitted notes. Hold out customers and periods; retain errors and
abstentions; examine whether the text precedes treatment or is itself a
consequence of contact. A measurement-error variance below $1/3$ is a boundary
in this particular constructed system, not a transferable acceptance threshold
for real text.

\section{Public banking records: an impressive score can use unavailable information}
The UCI Bank Marketing archive provides an actual public-data test of the
information-date requirement. The selected full additional file contains
41,188 archived contact records. Its documentation states that the rows are
ordered by date, from May 2008 to November 2010, and explicitly warns that
call duration is unavailable before the call. The smaller overlapping sample
is excluded from this experiment.
\citep{pd15_uci_bank}

The new study freezes the data identity, features, preprocessing, splits and
models before fitting. Its primary split trains on the first 28,831 rows and
tests on the remaining 12,357. A separately specified random stratified split
of the same sizes provides a development comparison, using seed 20261029.
Neither establishes independent customers: the archive lacks a suitable
customer identifier, and ordered rows do not supply exact decision timestamps.
No independent-observation confidence interval is claimed.

The baseline predicts the training subscription rate. The candidate model is
an L2-regularised logistic regression using age, job, marital status, education,
credit default, housing loan, personal loan, days since previous contact,
previous contacts and previous campaign outcome. Their pre-call availability
remains a maintained assumption requiring confirmation in a deployment. The
sentinel 999 for days since previous contact is represented by a separate
indicator and a zero numerical value; unknown categories remain categories.
Scaling and one-hot encoding are fitted on training data only. The deliberately
infeasible comparison adds completed call duration. Campaign count, current
contact fields and macroeconomic aggregates are excluded because their
decision-time meaning requires further clarification.

Both regressions use the same prespecified settings: $C=1$, the L-BFGS solver
and a maximum of 2,000 iterations. There is no tuning on test outcomes.
All four fits converge without recorded warnings, and no test categorical
values are absent from their corresponding training set.

\begin{table}[htbp]\centering\small
\begin{tabular}{llrrr}
\toprule
Split & Model & AUC & Log loss & Brier score \\
\midrule
Later archive & Training-rate constant & 0.500 & 0.752 & 0.221 \\
Later archive & Candidate pre-call & 0.630 & 0.678 & 0.206 \\
Later archive & Plus unavailable duration & 0.793 & 0.665 & 0.193 \\
Random development & Training-rate constant & 0.500 & 0.352 & 0.100 \\
Random development & Candidate pre-call & 0.693 & 0.312 & 0.087 \\
Random development & Plus unavailable duration & 0.883 & 0.254 & 0.073 \\
\bottomrule
\end{tabular}
\caption{Executed public-data prediction with a frozen specification. AUC
measures ranking; lower log loss and Brier score indicate better probability
predictions on the particular test set. Rows from different splits have
different outcome distributions and are not interchangeable performance tests.
The duration model is infeasible for the stated pre-call task.}
\label{scisyn:public-timing}
\end{table}

The apparent gain from duration is large: AUC rises from 0.630 to 0.793 on the
later archived records and from 0.693 to 0.883 on the random development split.
It cannot justify a pre-call ranking. Among the 1,236 highest-scored later
records, the candidate pre-call model identifies 620 recorded subscriptions,
while the duration model identifies 782. These are descriptive counts of
observed subscriptions among people already contacted, not estimates of extra
subscriptions caused by targeting them. The experiment supplies neither an
untreated comparison nor the assignment information needed for that causal
claim.

The ordered split also exposes a substantial composition change: the
subscription frequency is 5.57 percent in the early training rows and 24.55
percent in the later test rows. The random split has approximately 11.27 and
11.26 percent respectively. This contributes to the different prediction
problem; the experiment does not identify the cause of the change. AUC,
calibration and economic targeting value should not be collapsed into a single
claim that the model generalises.

The practical implication is observable without private bank data. A common
study interface should record which variables exist at the action date and
which records support the claimed evaluation. A simple rule excluding duration
already prevents one misleading improvement. Language assistance would need
to outperform that rule or uncover a harder omission to justify its extra
cost. The next experiment should compare completed analyst studies with these
requirements enforced, rather than rewarding the largest retrospective score.

